% !TeX program = xelatex
% 부산대학교 졸업과제 중간보고서 LaTeX 양식
% 컴파일 권장: XeLaTeX
% 부산대학교 로고 파일명: pusan_logo.png

\documentclass[12pt,a4paper,oneside]{report}

% -------------------------------------------------
% 기본 패키지
% -------------------------------------------------
\usepackage{kotex}
\usepackage{fontspec}
\usepackage{geometry}
\usepackage{graphicx}
\usepackage{float}
\usepackage{caption}
\usepackage{subcaption}
\usepackage{booktabs}
\usepackage{tabularx}
\usepackage{longtable}
\usepackage{array}
\usepackage{enumitem}
\usepackage{titlesec}
\usepackage{titletoc}
\usepackage{fancyhdr}
\usepackage{lastpage}
\usepackage[hidelinks]{hyperref}
\usepackage{xcolor}
\usepackage{setspace}
\usepackage{ragged2e}

% -------------------------------------------------
% 페이지 설정
% -------------------------------------------------
\geometry{
  top=22mm,
  bottom=20mm,
  left=22mm,
  right=22mm,
  headheight=14pt,
  footskip=10mm
}

\setstretch{1.35}
\setlength{\parindent}{1.5em}
\setlength{\parskip}{0pt}
\renewcommand{\arraystretch}{1.25}

% -------------------------------------------------
% 한글/영문 폰트
% 시스템에 없는 경우 아래 폰트명을 설치된 폰트로 바꾸세요.
% -------------------------------------------------
\IfFontExistsTF{Noto Sans CJK KR}{
  \setmainhangulfont{Noto Sans CJK KR}
  \setsanshangulfont{Noto Sans CJK KR}
}{
  \setmainhangulfont{UnDotum}
  \setsanshangulfont{UnDotum}
}

% -------------------------------------------------
% 보고서 기본 정보: 여기만 수정하면 표지에 반영됩니다.
% -------------------------------------------------
\newcommand{\ProjectTitle}{프로젝트 제목을 입력하세요}
\newcommand{\TeamName}{팀 이름}

\newcommand{\MemberOne}{202000000 홍길동}
\newcommand{\MemberTwo}{202000001 김철수}
\newcommand{\MemberThree}{202000002 이영희}
% 인원이 더 많으면 아래 표지 부분에 줄을 추가하세요.

\newcommand{\AdvisorName}{홍길동}
\newcommand{\SubmissionDate}{2026년 00월 00일}
\newcommand{\UniversityName}{부산대학교}

% -------------------------------------------------
% 장/절 제목 스타일
% 원본 보고서의 "제 1장 서론", "1.1 연구 배경..." 형식을 재현
% -------------------------------------------------
\titleformat{\chapter}[block]
  {\normalfont\bfseries\fontsize{24pt}{30pt}\selectfont}
  {제 \thechapter 장}
  {0.8em}
  {}
\titlespacing*{\chapter}{0pt}{0pt}{28pt}

\titleformat{\section}
  {\normalfont\bfseries\fontsize{17pt}{22pt}\selectfont}
  {\thesection}
  {0.7em}
  {}
\titlespacing*{\section}{0pt}{18pt}{10pt}

\titleformat{\subsection}
  {\normalfont\bfseries\fontsize{14pt}{19pt}\selectfont}
  {\thesubsection}
  {0.7em}
  {}
\titlespacing*{\subsection}{0pt}{14pt}{7pt}

% -------------------------------------------------
% 목차 스타일
% -------------------------------------------------
\renewcommand{\contentsname}{목차}
\titlecontents{chapter}
  [0pt]
  {\vspace{0.8em}\bfseries}
  {\contentslabel{2.3em}}
  {}
  {\titlerule*[0.7pc]{.}\contentspage}

\titlecontents{section}
  [1.5em]
  {}
  {\contentslabel{3.0em}}
  {}
  {}

\titlecontents{subsection}
  [3.0em]
  {}
  {\contentslabel{3.8em}}
  {}
  {}

% -------------------------------------------------
% 페이지 번호: 1/12 형식
% 표지와 목차에서는 숨기고 본문부터 표시
% -------------------------------------------------
\pagestyle{fancy}
\fancyhf{}
\renewcommand{\headrulewidth}{0pt}
\renewcommand{\footrulewidth}{0pt}
\fancyfoot[C]{\thepage/\pageref*{LastPage}}

% -------------------------------------------------
% 그림/표 캡션 형식: <그림 1. 설명>, <표 1. 설명>
% -------------------------------------------------
\renewcommand{\figurename}{그림}
\renewcommand{\tablename}{표}
\captionsetup{
  font=normalsize,
  labelfont=normalfont,
  labelsep=period,
  justification=centering,
  singlelinecheck=false
}

% 원본 스타일처럼 꺾쇠 캡션을 쓰고 싶을 때 사용
\newcommand{\figcaption}[1]{%
  \caption*{<그림 \thefigure. #1>}%
}
\newcommand{\tabcaption}[1]{%
  \caption*{<표 \thetable. #1>}%
}

% -------------------------------------------------
% 자주 쓰는 양식용 명령
% -------------------------------------------------
\newcommand{\placeholdertext}{%
  본문 내용을 입력하세요. 프로젝트의 목적, 구현 내용, 실험 결과 및 분석을 구체적으로 작성합니다.
}

\newcommand{\placeholderfigure}[2][0.85\textwidth]{%
  \begin{figure}[H]
    \centering
    \fbox{\parbox[c][6cm][c]{#1}{\centering 이미지 삽입 영역\\[0.5em]\texttt{\textbackslash includegraphics\{파일명.png\}}}}
    \caption{#2}
  \end{figure}
}

\newcommand{\placeholderwidefigure}[2][0.95\textwidth]{%
  \begin{figure}[H]
    \centering
    \fbox{\parbox[c][9cm][c]{#1}{\centering 큰 이미지/시스템 구조도 삽입 영역}}
    \caption{#2}
  \end{figure}
}

% -------------------------------------------------
% 문서 시작
% -------------------------------------------------
\begin{document}

% =================================================
% 표지
% =================================================
\begin{titlepage}
  \thispagestyle{empty}
  \begin{center}
    \vspace*{22mm}

    {\fontsize{25pt}{34pt}\selectfont\bfseries \ProjectTitle\par}

    \vspace{24mm}

    {\fontsize{16pt}{22pt}\selectfont\bfseries 팀 \TeamName\par}

    \vspace{10mm}

    {\fontsize{14pt}{22pt}\selectfont
      \MemberOne\\
      \MemberTwo\\
      \MemberThree
    \par}

    \vspace{16mm}

    {\fontsize{14pt}{20pt}\selectfont 지도교수 \AdvisorName\par}

    \vspace{12mm}

    {\fontsize{14pt}{20pt}\selectfont 제출일 \SubmissionDate\par}

    \vfill

    % 부산대학교 로고 파일을 같은 폴더에 pusan_logo.png 이름으로 두세요.
    \IfFileExists{pnulogo.jpg}{%
      \includegraphics[width=0.22\textwidth]{pusan_logo.png}\\[6mm]
    }{%
      \fbox{\parbox[c][28mm][c]{0.22\textwidth}{\centering 부산대학교 로고\\\texttt{pusan\_logo.png}}}\\[6mm]
    }

    {\fontsize{17pt}{22pt}\selectfont\bfseries \UniversityName\par}
    \vspace{10mm}
  \end{center}
\end{titlepage}

% =================================================
% 목차
% =================================================
\pagenumbering{gobble}
\tableofcontents
\clearpage

% =================================================
% 본문 시작
% =================================================
\pagenumbering{arabic}
\setcounter{page}{1}

% =================================================
% 제 1장 서론
% =================================================
\chapter{서론}

\section{연구 배경 및 기존 문제점}
\placeholdertext

\placeholderfigure{연구 배경 관련 그림 제목}

\begin{itemize}[leftmargin=1.5em,label={-},itemsep=2pt,topsep=4pt]
  \item 문제점 또는 배경 요약 1
  \item 문제점 또는 배경 요약 2
  \item 문제점 또는 배경 요약 3
\end{itemize}

\section{소프트웨어 설계 목표}
\placeholdertext

\placeholderfigure{페르소나 또는 설계 목표 관련 그림 제목}

\begin{itemize}[leftmargin=1.5em,label={-},itemsep=2pt,topsep=4pt]
  \item 설계 목표 1
  \item 설계 목표 2
  \item 설계 목표 3
\end{itemize}

% =================================================
% 제 2장 프로젝트 기능
% =================================================
\chapter{프로젝트 기능}

\section{주요 기능 소개}

\subsection{기능 1}
\placeholdertext

\subsection{기능 2}
\placeholdertext

\subsection{기능 3}
\placeholdertext

\subsection{기능 4}
\placeholdertext

\subsection{기능 5}
\placeholdertext

\subsection{기능 6}
\placeholdertext

\subsection{기능 7}
\placeholdertext

% 기능 화면 예시
\placeholderwidefigure{주요 기능 화면 1}
\placeholderwidefigure{주요 기능 화면 2}
\placeholderwidefigure{주요 기능 화면 3}

% =================================================
% 제 3장 시스템 설계 및 구현
% =================================================
\chapter{시스템 설계 및 구현}

\section{전체 구조 개요}

\subsection{클라이언트--서버 아키텍처}
\begin{itemize}[leftmargin=1.5em,label={-},itemsep=2pt,topsep=4pt]
  \item 클라이언트: 기술 스택 입력
  \item 서버: 기술 스택 입력
  \item 데이터베이스: 사용 기술 입력
  \item 외부 API 및 AI 기능: 사용 기술 입력
\end{itemize}

\placeholderwidefigure{서비스 아키텍처}

\subsection{주요 통신 흐름}
\begin{itemize}[leftmargin=1.5em,label={-},itemsep=2pt,topsep=4pt]
  \item 사용자 $\rightarrow$ 앱: 입력 또는 요청
  \item 앱 $\rightarrow$ 서버: API 요청
  \item 서버 $\rightarrow$ 외부 API: 데이터 조회 또는 AI 호출
  \item 서버 $\leftrightarrow$ DB: 데이터 저장 및 조회
\end{itemize}

\placeholderwidefigure{주요 처리 흐름}

\section{프론트엔드}

\subsection{기술 스택}
\begin{itemize}[leftmargin=1.5em,label={-},itemsep=2pt,topsep=4pt]
  \item 기술 1
  \item 기술 2
  \item 기술 3
\end{itemize}

\subsection{주요 화면 구조}
\begin{itemize}[leftmargin=1.5em,label={-},itemsep=2pt,topsep=4pt]
  \item 화면 1: 설명
  \item 화면 2: 설명
  \item 화면 3: 설명
\end{itemize}

\section{백엔드}

\subsection{기술 스택}
\begin{itemize}[leftmargin=1.5em,label={-},itemsep=2pt,topsep=4pt]
  \item 기술 1
  \item 기술 2
  \item 기술 3
\end{itemize}

\subsection{DB 설계 및 주요 테이블}
\begin{table}[H]
  \centering
  \begin{tabularx}{0.92\textwidth}{>{\bfseries}l X}
    \toprule
    테이블명 & 역할 \\
    \midrule
    User & 사용자 정보 \\
    Example & 주요 데이터 설명 \\
    \bottomrule
  \end{tabularx}
  \caption{주요 데이터베이스 테이블}
\end{table}

\subsection{보안 및 인증}
\placeholdertext

\subsection{AI 연동 로직}
\placeholdertext

% =================================================
% 제 4장 개발 일정 및 역할 분담
% =================================================
\chapter{개발 일정 및 역할 분담}

\section{개발 일정}

\subsection{1차 기획 및 설계}
\begin{itemize}[leftmargin=1.5em,label={-},itemsep=2pt,topsep=4pt]
  \item 아이디어 구체화 및 서비스 기획
  \item 기술 스택 선정 및 기본 설계
\end{itemize}

\subsection{2차 개발 및 통신 구현}
\begin{itemize}[leftmargin=1.5em,label={-},itemsep=2pt,topsep=4pt]
  \item 프론트엔드 구현
  \item 백엔드 및 API 연동
\end{itemize}

\subsection{3차 테스트 및 마무리}
\begin{itemize}[leftmargin=1.5em,label={-},itemsep=2pt,topsep=4pt]
  \item 주요 기능 통합 테스트
  \item 발표 자료 및 보고서 작성
\end{itemize}

\section{역할 분담}

\subsection{팀원 1}
\begin{itemize}[leftmargin=1.5em,label={-},itemsep=2pt,topsep=4pt]
  \item 담당 업무 1
  \item 담당 업무 2
\end{itemize}

\subsection{팀원 2}
\begin{itemize}[leftmargin=1.5em,label={-},itemsep=2pt,topsep=4pt]
  \item 담당 업무 1
  \item 담당 업무 2
\end{itemize}

\subsection{팀원 3}
\begin{itemize}[leftmargin=1.5em,label={-},itemsep=2pt,topsep=4pt]
  \item 담당 업무 1
  \item 담당 업무 2
\end{itemize}

% =================================================
% 제 5장 결과 분석 및 기대 효과
% =================================================
\chapter{결과 분석 및 기대 효과}

\section{결과 분석}

\subsection{분석 항목 1}
\placeholdertext

\subsection{분석 항목 2}
\placeholdertext

\subsection{분석 항목 3}
\placeholdertext

\subsection{시스템 안정성 및 확장성}
\placeholdertext

\section{기대 효과}
\begin{itemize}[leftmargin=1.5em,label={-},itemsep=2pt,topsep=4pt]
  \item 기대 효과 1
  \item 기대 효과 2
  \item 기대 효과 3
\end{itemize}

% =================================================
% 제 6장 결론 및 향후 계획
% =================================================
\chapter{결론 및 향후 계획}

\section{결론}
\placeholdertext

\section{향후 개발 계획}

\subsection{고도화 방향}
\placeholdertext

\subsection{추후 플랫폼 확장}
\placeholdertext

\section{기대 성장 방향}
\placeholdertext

% =================================================
% 제 7장 참고문헌
% =================================================
\chapter{참고문헌}

\begin{enumerate}[leftmargin=2em,label=\arabic*.]
  \item 저자, ``문헌 제목'', 출판사 또는 학회명, 연도.
  \item 공식 문서명, 서비스 또는 기관명, \url{https://example.com}.
  \item 저자, ``논문 제목'', 학회/저널명, 권(호), 페이지, 연도.
\end{enumerate}

\end{document}
